\section{Introduction}
\label{sec:introduction}

Web-scale text corpora are increasingly contaminated by synthetic content that has itself been generated from earlier synthetic content.
Estimates suggest that a substantial fraction of online text is machine-generated~\cite{sun2025aigt, spennemann2025quantifying}, and this share continues to rise.
When models train on the outputs of earlier models, the resulting multi-generational synthetic text degrades with each iteration: tail modes vanish, lexical diversity contracts, and surprisal distributions flatten~\cite{shumailov2024collapse, dohmatob2024model}.
This degradation arises in both \emph{retrain-chains} (iteratively training on model outputs) and \emph{prompt-chains} (recursively generating from previous outputs without retraining); we study the latter, which isolates generation-level information loss from training dynamics (Appendix~\ref{sec:extended_limitations}).
Yet the dominant paradigm treats all machine-generated text identically, discarding the generational depth dimension.

This gap cuts across four lines of work: binary detectors~\cite{basani2026diveye, zhang2024minkpp} produce no depth estimate (Section~\ref{sec:binary_confidence}); membership inference~\cite{zawalski2026codec} addresses a different question; the model collapse literature~\cite{shumailov2024collapse, alemohammad2024selfconsuming, schaeffer2025position} provides no per-text tool; and data selection methods~\cite{xie2023dsir, tirumala2023d4, drayson2025} ignore recursive depth.
No existing method estimates how many rounds of recursive generation a text has undergone.

The Data Processing Inequality guarantees that information degrades monotonically across generations, motivating a finite discrimination boundary $\Kstar$.
We empirically measure a lower bound $\Kempirical$ on the discrimination boundary (the maximum depth at which a classifier reliably distinguishes adjacent generations) using 15 statistical features and an Ordinal Binary Decomposition (OBD) LightGBM classifier.
We operationalize $\Kempirical$ as the deepest $k \geq 1$ where the pairwise area under the ROC curve ($\auc_{k-1,k}$) exceeds $\theta = 0.60$ (Section~\ref{sec:formulation}).

Our contributions:
\begin{enumerate}[label=(\alph*),nosep]
  \item \textbf{Measuring the discrimination boundary}: $\Kempirical \leq 2$ across all 10 model configurations (6 architecture families at 1B--8B scale) and 4 decoding strategies; $\Kempirical{=}2$ for 6/10 spanning three families ($p < 0.001$; Section~\ref{sec:kstar_boundary}).
  Beyond two generation passes, web text becomes statistically indistinguishable across deeper generations; this finding establishes binary sufficiency by measurement, reveals domain-dependent variation (arXiv $\Kempirical{=}3$), and quantifies boundary sensitivity to threshold, scorer, and scale.
  Both supervised OBD and zero-shot large language models (LLMs) fail beyond $\Kempirical$ (McNemar $p > 0.05$; Section~\ref{sec:why_kstar}).
  \item \textbf{Domain dependence of $\Kempirical$}: general web text (C4, Wikipedia) and code yield $\Kempirical{\leq}2$, while structured academic text (arXiv) yields $\Kempirical{=}3$, consistent across model families (Section~\ref{sec:cross_domain}).
  \item \textbf{Generator-agnostic cross-model transfer} (within continuation mode): 15 statistical features enable 80--92\% accuracy retention with scorer-invariant feature rankings ($\rho > 0.90$), supporting depth estimation on corpora from unknown continuation-mode generators without retraining; $\Kempirical$ agreement holds across scorers for strong-signal configurations (Section~\ref{sec:transfer}).
\end{enumerate}

\noindent As a consistency check under extreme contamination (${\sim}$83\% synthetic), $\Kempirical{=}2$ correctly predicts the strategy ordering: binary curation minimizes Massive Multitask Language Understanding (MMLU) degradation ($0.5941$ vs.\ base $0.6018$); at realistic ratios ($\leq$50\%), all strategies converge (Section~\ref{sec:filtering}).
